\section{Một tiên đề ở tầng trên}
\label{sec:ch11-tien-de-o-tang-tren}

\index{tính hiệu quả phân tầng|(}%
Một xây dựng như thế được bảo đảm bằng cái gì? Bài trả lời bằng một tiên đề mà
người đọc Phần~II đã gặp hình dạng của nó hai lần rồi, và đặt cho nó cái tên
\tn{tính hiệu quả phân tầng}{hierarchical efficiency}.

\begin{definition}[Tính hiệu quả phân tầng]
  \label{def:ch11-hieu-qua-phan-tang}
  Một họ đại lượng bậc hai $\varphi_{j \to i}$ được gọi là hiệu quả phân tầng
  đối với attribution bậc nhất $\varphi_i(x,f)$ nếu
  %
  \begin{equation}
    \label{eq:ch11-hieu-qua-phan-tang}
    \varphi_i(x,f)
      = \sum_{j \in [d]} \varphi_{j \to i}(x,f)
      = \varphi_{i \to i}(x,f) + \sum_{j \neq i} \varphi_{j \to i}(x,f)
  \end{equation}
  %
  với mọi đặc trưng $i$~\autocite{p23metagame}.
\end{definition}

Phương trình~\ref{eq:ch11-hieu-qua-phan-tang} nói rằng attribution bậc nhất của
một đặc trưng được lấy lại nguyên vẹn khi cộng \tn{hiệu ứng riêng}{pure individual effect} của nó với các
hiệu ứng tương tác của nó với những đặc trưng còn lại. Không có phần nào của
$\varphi_i$ rơi ra ngoài, và không có phần nào được đếm hai lần. Bài chứng minh
rằng meta-attribution của định nghĩa~\ref{def:ch11-meta-attribution} thỏa tính
chất ấy, nên phép phân tách mà mục trước dựng không đánh mất gì của tầng dưới.

Hình dạng ấy chính là hình dạng của \tn{tính đầy đủ}{completeness}, tiên đề mà
mệnh đề~\ref{prop:ch05-tinh-day-du} phát biểu cho Integrated Gradients, và của
\tn{độ chính xác cục bộ}{local accuracy} mà chương~\ref{ch:shap} rút ra từ tiên đề của giá trị
Shapley. Cả ba nói cùng một điều ở ba tầng khác nhau: tổng các phần bằng cái
được chia. Ở tính đầy đủ, cái được chia là hiệu số đầu ra của mô hình giữa đầu
vào và baseline, còn các phần là attribution của từng đặc trưng. Ở tính hiệu
quả phân tầng, cái được chia là attribution của một đặc trưng, còn các phần là
meta-attribution của đặc trưng ấy với từng đặc trưng khác. Tiên đề đi lên một
tầng cùng với xây dựng.

Kết quả thứ nhất của bài nói rằng tiên đề ấy đã được thỏa từ trước, ở những chỗ
không ai phát biểu nó.

\begin{theorem}[Các chỉ số tương tác sẵn có đều hiệu quả phân tầng]
  \label{thm:ch11-phan-tang-san-co}
  Các chỉ số tương tác đang được dùng đều thỏa
  phương trình~\ref{eq:ch11-hieu-qua-phan-tang}. Cả họ tương tác Shapley, mà
  STII là một thành viên, phân tách giá trị Shapley, còn SOP phân tách
  Integrated Gradients, với quy ước rằng
  đại lượng bậc hai lấy từ một chỉ số theo tập là giá trị của chỉ số ấy trên
  cặp chia cho số phần tử của cặp. Các phương pháp nối tầng cũng phân tách
  attribution bậc nhất của chúng, với đại lượng bậc hai lấy thẳng bằng giá trị
  nối tầng~\autocite{p23metagame}.
\end{theorem}

Bài gọi kết quả này là điều đáng ngạc nhiên, vì những chỉ số ấy được xây dựng
với mục đích khác và không ai trình bày chúng như một phép phân tách theo tầng.
Kết quả thứ hai làm chặt thêm quan hệ: Meta-SV và Meta-IG là biến thể có hướng
của hai chỉ số theo tập đã có, cụ thể là của STII và của SOP. Với mỗi cặp
$i \neq j$,
%
\begin{equation}
  \label{eq:ch11-doi-xung-hoa}
  \text{STII}_{\{i,j\}}
    = \varphi^{\text{Meta-SV}}_{j \to i} + \varphi^{\text{Meta-SV}}_{i \to j},
  \qquad
  \text{SOP}_{\{i,j\}}
    = \varphi^{\text{Meta-IG}}_{j \to i} + \varphi^{\text{Meta-IG}}_{i \to j},
\end{equation}
%
nghĩa là chỉ số theo tập thu được từ đại lượng có hướng bằng cách cộng hai
chiều lại~\autocite{p23metagame}. Đọc theo chiều ngược, đó là lời giải thích
cho ô gộp trong bảng~\ref{tab:ch11-tach-hieu-ung}: chỉ số theo tập không mất
thông tin về chiều vì nó chưa bao giờ có, mà vì nó cộng hai chiều lại trước khi
báo cáo.

Đến đây phải tách ra điều mà tiên đề bảo đảm với điều nó không bảo đảm, và chỗ
tách này là cách đọc của sách chứ không phải phát biểu của bài. Tính hiệu quả
phân tầng ràng buộc quan hệ giữa tầng hai và tầng một. Nó không nói gì về việc
tầng một có đúng hay không, vì $\varphi_i$ đứng ở vế trái của
phương trình~\ref{eq:ch11-hieu-qua-phan-tang} như một đại lượng cho trước.

Gradient nhân đầu vào là minh họa sẵn có cho điều đó, và bài in các con số cần
thiết ở hai cột mà bảng~\ref{tab:ch11-tach-hieu-ung} không chép lại. Với hàm ở
phương trình~\ref{eq:ch11-ham-vi-du}, phương pháp ấy cho
$\varphi_1 = x_1 + T$ và $\varphi_2 = 2T$, nên tổng hai attribution bậc nhất
của nó bằng $x_1 + 3T$, trong khi hiệu số đầu ra giữa đầu vào và baseline chỉ
bằng $x_1 + T$~\autocite{p23metagame}. Phương pháp bậc nhất ấy vi phạm tính đầy
đủ, và Meta-G$\times$I dựng trên nó vẫn hiệu quả phân tầng: cộng hàng
Meta-G$\times$I của bảng lại thì được đúng $\varphi_1$ và đúng $\varphi_2$, tức
là phép phân tách bảo toàn nguyên vẹn một con số vốn đã sai.

Đó là kết luận mà chương~\ref{ch:attribution-tu-nguyen-ly-dau} đã rút ra ở tầng
dưới, gặp lại ở tầng trên. Mục~\ref{sec:ch06-do-do-la-gia-thuyet} cho thấy các
tiên đề của attribution ràng buộc được dạng của lời giải thích mà không chọn
được \tn{độ đo}{measure}, nên một xây dựng thỏa hết tiên đề vẫn còn nguyên câu hỏi nó có đúng
hay không. Ở đây cũng vậy, chỉ khác là câu hỏi ấy giờ có hai tầng để hỏi.
Mục~\ref{sec:ch11-ba-phep-do-va-ba-cai-nen} xem bài trả lời nó bằng cách nào.
\index{tính hiệu quả phân tầng|)}%
